% ============================================================
% V. DISCUSIÓN
% ============================================================

\section{Discusión}

Esta sección responde las preguntas de reflexión planteadas por la guía de práctica. Al no disponerse de datos reales (Sección~\ref{sec:resultados}), las respuestas se fundamentan en el comportamiento teórico esperado del sistema implementado, señalando explícitamente qué debería verificarse una vez existan capturas reales.

\subsection{Digitalización de señales periódicas}

\textbf{Causas de las diferencias entre frecuencia normalizada calculada y observada.} Aun con una implementación correcta, se esperan diferencias no nulas por varias razones independientes: (1) la resolución en frecuencia de la FFT es $f_s/N$, por lo que el pico observado solo puede ubicarse en un múltiplo discreto de esa resolución, nunca exactamente en el valor teórico salvo coincidencia; (2) el error máximo permitido del $5\,\%$ en el periodo de la interrupción de temporizador (Sección de Metodología) se traduce directamente en un error proporcional de la frecuencia de muestreo efectiva $f_s$ y, por tanto, de toda frecuencia normalizada calculada con base en el valor nominal de $500$ Hz; (3) el generador de ondas de laboratorio tiene su propia tolerancia de frecuencia, no infinitamente precisa; y (4) cuando la ventana de captura no contiene un número entero de ciclos de la señal —como ocurre con $75$ Hz, que produce $6.667$ muestras/ciclo, no entero, según la Tabla~\ref{tab:periodica}—, se produce \emph{fuga espectral} (\emph{spectral leakage}): la energía de la componente sinusoidal se dispersa entre varios bins de la FFT en lugar de concentrarse en uno solo, ensanchando el pico observado y desplazando su máximo respecto a la frecuencia real.

\textbf{¿Están presentes en todas las filas?} Se espera que sí, en algún grado, dado que las causas (1)–(3) son inherentes al sistema y no dependen de la frecuencia específica. Sin embargo, la magnitud de la diferencia no debería ser uniforme: se espera que sea mayor para $75$ Hz que para $50$, $100$ y $125$ Hz, precisamente por la fuga espectral adicional de la causa (4), ya que estas tres últimas sí producen un número entero de muestras por ciclo dentro de una ventana de captura cuya duración sea un múltiplo entero del periodo de la señal.

\subsection{Análisis frecuencial}

El pico en $0$ Hz del espectro corresponde al nivel de continua (DC) de la señal, y su altura es directamente proporcional a la magnitud de dicho nivel DC —de hecho, tras la normalización $2/N$ empleada (Sección~\ref{sec:muestreo}), la altura del pico en $0$ Hz converge exactamente al valor del nivel DC configurado en el generador. La ubicación del pico correspondiente a la componente sinusoidal coincide con la frecuencia de dicha componente, y su amplitud converge al valor pico de la sinusoide (o a la mitad de su amplitud pico-a-pico) bajo la misma normalización. La formulación de ajuste sugerida es exactamente la implementada en \texttt{matlab/analyze\_periodic\_signal.m}: graficar $\text{mag} = \frac{2}{N}\left|X[k]\right|$ contra $\text{freq} = k \cdot f_s/N$ para $k = 0,\dots,N/2$, de modo que el eje horizontal quede directamente en Hz y el eje vertical directamente en las unidades físicas de la señal (voltios), sin factores de escala adicionales que el lector deba inferir.

\subsection{Reconstrucción de la señal digitalizada}

\textbf{Relación entre muestras por ciclo y escalones de la reconstrucción DAC.} Como se explica en la Sección~\ref{sec:reconstruccion}, cada muestra capturada produce exactamente un escalón en la salida del DAC (retenedor de orden cero), de modo que el número de escalones visibles por ciclo en la señal reconstruida es, por construcción, idéntico al número de muestras por ciclo calculado en la Tabla~\ref{tab:periodica} para esa misma frecuencia —no es una coincidencia experimental, sino una consecuencia directa de que el DAC solo puede cambiar de valor una vez por interrupción de muestreo.

\textbf{Causas del desfase de la reconstrucción PWM + filtro.} El desfase observado se debe principalmente al retardo de grupo introducido por el filtro pasa-bajos analógico externo, que depende de su orden y de su frecuencia de corte respecto a la componente de la señal; a este se suma un retardo menor y común a ambos métodos, el del propio retenedor de orden cero del muestreo (en promedio, medio periodo de muestreo). \textbf{¿Son iguales los desfases de PWM y DAC?} No necesariamente: el método DAC solo hereda el retardo del retenedor de orden cero, mientras que el método PWM añade, además, el retardo de grupo del filtro externo, que típicamente lo domina y por tanto produce un desfase total distinto (y usualmente mayor) que el del DAC directo. \textbf{¿Son los desfases constantes con la frecuencia de entrada?} El retardo del retenedor de orden cero es aproximadamente constante en muestras (medio periodo de muestreo), por lo que en tiempo absoluto no depende de la frecuencia de la señal de entrada; el retardo de grupo del filtro analógico, en cambio, generalmente sí varía con la frecuencia salvo que el filtro se diseñe específicamente con fase lineal en la banda de interés — por lo que el desfase del método PWM no debería considerarse constante en general.

\textbf{Ventajas, desventajas y aplicaciones.} El método DAC directo es más simple de implementar en firmware, no requiere circuitería analógica adicional, y su desfase es más predecible (depende solo del muestreo, no de un filtro externo); su desventaja es que su salida es intrínsecamente escalonada, con contenido espectral de alta frecuencia (armónicos del muestreo) que puede requerir su propio filtrado si se necesita una salida verdaderamente suave. El método PWM + filtro produce, tras un filtro bien diseñado, una salida más suave sin depender de la resolución del DAC de la tarjeta (útil si esta no tiene DAC dedicado, o si se necesita mayor potencia de salida mediante una etapa de conmutación); su desventaja es la complejidad y el costo adicional del filtro analógico externo, su retardo de grupo generalmente mayor y dependiente de la frecuencia, y la necesidad de un diseño de filtro cuidadoso para lograr suficiente atenuación de la portadora PWM sin distorsionar la señal de interés. En términos de aplicación, el DAC directo es preferible cuando se requiere baja latencia y simplicidad (por ejemplo, generación de referencias de control), mientras que PWM + filtro es preferible en escenarios de potencia (control de motores, actuadores) donde de todas formas se dispone de una etapa de conmutación y el filtrado forma parte natural del sistema.

\subsection{Digitalización de sensores — BME280}

La guía de práctica original formula esta subsección en torno a una IMU; al usarse un BME280 (presión, temperatura, humedad) en su lugar, no existen canales de aceleración ni de velocidad angular, y la comparación ``quieto'' vs. ``en movimiento'' no aplica directamente a un sensor sin partes ni ejes inerciales. Las preguntas se responden a continuación adaptadas a la comparación real implementada: condiciones \textbf{base} (ambiente estable) vs. \textbf{perturbadas} (Sección~\ref{sec:resultados}).

\textbf{Diferencias entre condiciones base y perturbadas.} En condiciones base, con el sensor en un ambiente térmicamente estable y sin fuentes cercanas de calor, humedad o presión diferencial, se espera que los tres canales muestren una media aproximadamente constante (la temperatura y humedad ambientales, y la presión atmosférica local) con una desviación estándar pequeña, atribuible principalmente al ruido electrónico interno del sensor y a la resolución efectiva de las fórmulas de compensación. Bajo perturbación deliberada (por ejemplo, una fuente de calor o el aliento acercados al sensor, o cubrirlo parcialmente), se espera un cambio significativo en la media de temperatura y humedad —ambas suben con el aliento o una fuente de calor cercana—, y un cambio más sutil, o incluso nulo, en la presión, ya que esta responde principalmente a la altitud/condiciones atmosféricas y no a perturbaciones locales pequeñas salvo que se module directamente el volumen de aire alrededor del sensor (por ejemplo, cubriéndolo).

\textbf{Identificación y características del ruido.} En condiciones base, el ruido de los tres canales es de banda ancha y de amplitud pequeña relativa al nivel DC dominante (temperatura/humedad/presión ambientales), consistente con ruido térmico/electrónico propio del sensor y con el redondeo introducido por la compensación en punto flotante —su ancho de banda calculado por \texttt{matlab/analyze\_bme280.m} debería ser relativamente amplio pero de magnitud pequeña. Bajo perturbación, el cambio dominante es un transitorio de baja frecuencia (el tiempo de establecimiento térmico/higrométrico del propio sensor, típicamente de segundos, mucho más lento que los $100$ Hz de muestreo) superpuesto al mismo ruido de fondo de banda ancha; a diferencia de un sensor inercial, aquí no se espera contenido de alta frecuencia genuino asociado a la perturbación, ya que las magnitudes físicas medidas (temperatura, presión, humedad) son intrínsecamente de dinámica lenta.

\textbf{Estrategias de mitigación de ruido.} Las estrategias aplicables incluyen filtrado digital pasa-bajos (un filtro de media móvil es especialmente adecuado, dada la dinámica lenta de las tres magnitudes), promediado de múltiples muestras —tolerable aquí sin penalizar la respuesta a eventos de interés, ya que estos ocurren en la escala de segundos, no de milisegundos—, aumento del sobremuestreo interno del propio BME280 (\texttt{ctrl\_meas}/\texttt{ctrl\_hum}, a costa de una tasa de muestreo aún menor que los $100$ Hz ya limitados por el sensor), y calibración de \emph{offset} en condiciones ambientales conocidas antes de cada sesión de medición si se requiere precisión absoluta, no solo relativa.

\subsection{Digitalización de sensores — encoder y motorreductor}

\textbf{Relación velocidad/voltaje aplicado.} Para un motor DC en estado estacionario y bajo carga aproximadamente constante, la relación entre la velocidad angular y el voltaje aplicado es aproximadamente lineal, consistente con el modelo de estado estable de un motor DC $\omega \approx (V - I R_a)/K_e$, donde para una carga fija la caída resistiva $I R_a$ varía poco frente a $V$, dejando una relación cercana a $\omega \propto V$ con un posible desplazamiento (voltaje mínimo de arranque, por fricción estática) antes de que el motor comience a girar.

\textbf{Características del ruido.} La señal de velocidad angular derivada del conteo de pulsos por intervalo tiene una fuente de ruido intrínseca de cuantización: el número de pulsos contados en cada ventana de $5$ ms es un entero, por lo que la velocidad estimada tiene una resolución mínima de $\frac{2\pi}{\texttt{PULSES\_PER\_REV} \times 0.005\,\text{s}}$ rad/s, que se manifiesta como ruido de cuantización de alta frecuencia superpuesto a la velocidad real, además de la variación genuina de velocidad instantánea del motor bajo carga.

\textbf{Estrategias de reducción de ruido.} Aumentar el tamaño de la ventana de conteo reduce el ruido de cuantización relativo (a costa de menor resolución temporal), promediar velocidades de varias ventanas consecutivas, o cambiar a la estrategia de medición de tiempo entre pulsos a velocidades bajas (donde produce menor ruido relativo que el conteo por intervalo) son las estrategias más directas; un filtro pasa-bajos digital adicional sobre la señal de velocidad estimada también es aplicable si la dinámica de interés es mucho más lenta que la resolución del encoder lo permite.

\subsection{Fenómeno de aliasing}
\label{sec:discusion-aliasing}

\textbf{¿Qué ocurre con los datos capturados cuando $f \geq 250$ Hz?} Al superar la frecuencia de Nyquist ($f_s/2 = 250$ Hz), la señal capturada deja de representar fielmente la frecuencia real del generador: en su lugar, el sistema captura una señal cuya frecuencia normalizada observada coincide con la frecuencia \emph{plegada} teórica de la Tabla~\ref{tab:aliasing} (Ec.~\ref{eq:aliasing}), indistinguible —a partir únicamente de los datos digitalizados— de una señal genuina de esa frecuencia más baja. Es exactamente lo que predice la columna de frecuencia normalizada calculada de la Tabla~\ref{tab:aliasing}: $550$, $950$, $1050$ y $2050$ Hz deberían observarse todos como $0.100$ ciclos/muestra ($50$ Hz aparentes); $575$, $1075$ y $2075$ Hz como $0.150$ ($75$ Hz aparentes); y $500$, $1000$ y $2000$ Hz como $0.000$, es decir, indistinguibles de una señal de continua (DC) pura.

\textbf{Conclusión sobre el proceso de digitalización.} La coincidencia teórica de frecuencias normalizadas entre múltiplos de $f_s$ separados exactamente $500$ Hz confirma que el proceso de muestreo es fundamentalmente periódico en el dominio de la frecuencia: cualquier frecuencia analógica y sus reflejos alrededor de los múltiplos de $f_s/2$ son indistinguibles una vez digitalizados, sin importar cuán distintas sean antes de muestrear.

\textbf{Explicación del fenómeno de aliasing a partir de los resultados.} El aliasing ocurre porque el proceso de muestreo, al tomar solo muestras discretas y periódicas de la señal, no conserva suficiente información para distinguir entre una componente de frecuencia $f$ y su reflejo $|f - kf_s|$ dentro de la banda $[0, f_s/2]$: ambas producen exactamente la misma secuencia de muestras. La reconstrucción posterior (sea DAC o PWM+filtro) no puede recuperar cuál de las dos frecuencias originó realmente los datos, y por tanto reproduce la frecuencia plegada, no la real —de ahí que la señal reconstruida en el osciloscopio para $f \geq 250$ Hz debería mostrar la frecuencia aparente de la Tabla~\ref{tab:aliasing}, no la frecuencia realmente configurada en el generador.

\textbf{Máxima frecuencia digitalizable sin afectar la información.} De acuerdo con el teorema de Nyquist, la máxima frecuencia que puede digitalizarse sin ambigüedad con $f_s = 500$ Hz es $f_s/2 = 250$ Hz —el propio límite en el que la Tabla~\ref{tab:aliasing} comienza su primera fila—, ya que cualquier frecuencia por encima de ese límite queda indistinguible de alguna frecuencia dentro de $[0, 250]$ Hz una vez muestreada.
